\documentclass[11pt,a4paper]{altacv}

\geometry{left=1.5cm,right=1.5cm,top=1.cm,bottom=1.2cm,columnsep=1.5cm}

\usepackage[utf8]{inputenc}
\usepackage[T1]{fontenc}
\usepackage[french]{babel}
\usepackage{paracol}
\usepackage{xcolor}
\usepackage{hyperref}
\usepackage{fontawesome5}
\usepackage{setspace}

\definecolor{SlateGrey}{HTML}{2E2E2E}
\definecolor{LightGrey}{HTML}{666666}
\definecolor{DarkPastelRed}{HTML}{8AC0D9}
\definecolor{PastelRed}{HTML}{00B0DA}
\definecolor{GoldenEarth}{HTML}{B4AD0C}
\colorlet{name}{black}
\colorlet{tagline}{PastelRed}
\colorlet{heading}{DarkPastelRed}
\colorlet{headingrule}{GoldenEarth}
\colorlet{subheading}{PastelRed}
\colorlet{accent}{PastelRed}
\colorlet{emphasis}{SlateGrey}
\colorlet{body}{LightGrey}

\renewcommand{\familydefault}{\sfdefault}

\name{\Large Guillaume BOILEAU}
\tagline{Chef de projet technique SI | SPOC, besoins techniques, build/run, risques, cycle en V \& coordination DSI}
\personalinfo{
  \email{guillaume.boileaupro@gmail.com}
  \phone{+33 6 59 94 85 84}
  \location{Alpes-Maritimes, France}
  \linkedin{guillaume-boileau-phd-891b81265}
}

\setlength{\parskip}{0.72\baselineskip}

\begin{document}
\makecvheader
\vspace{-0.65cm}

\cvsection{Lettre de motivation}
\vspace{-0.45cm}

{\color{accent}\textbf{Objet :}} \textit{Candidature -- Chef de projet stratégique informatique -- Valbonne}

{\color{body}

Madame, Monsieur,

Je vous adresse ma candidature pour le poste de \textbf{Chef de projet stratégique informatique} au sein de la Direction Technique de la CNAF. Ce qui m'intéresse dans votre offre, c'est le rôle très concret de point d'entrée technique : qualifier les demandes, clarifier les contraintes, consolider les impacts build/run, fédérer les acteurs DT/DO, suivre les actions, identifier les risques et préparer les arbitrages en comitologie.

Mon profil correspond à cette logique d'interface. Je viens d'environnements techniques complexes où la valeur ne consiste pas seulement à produire du logiciel, mais à rendre les sujets compréhensibles, pilotables et transférables : structurer une demande, identifier les dépendances, documenter les choix, suivre les écarts, alerter au bon moment et maintenir les parties prenantes au même niveau d'information.

Au \textbf{CNES}, j'ai travaillé sur le segment sol de la mission spatiale \textbf{LISA}, dans un contexte CNES/ESA combinant chaînes de données mission L0--L3, cycle en V, intégration de modèles, validation et revues techniques. J'ai développé \textbf{CATpyLISA}, une plateforme Python utilisée pour structurer des workflows, comparer des résultats, qualifier des écarts et produire une documentation exploitable par plusieurs contributeurs. Cette expérience m'a donné une pratique solide du cadrage technique, de la validation, de la traçabilité, de la synthèse et du transfert de connaissances.

Chez \textbf{VEV Platform Services France}, j'ai complété ce socle par une expérience plus opérationnelle sur une plateforme logicielle connectée à des infrastructures terrain. J'y ai travaillé sur l'intégration de services, les flux de données, l'analyse d'incidents, le support run et la formalisation de constats techniques. Cette expérience est directement utile pour comprendre les enjeux de production, de maintien en condition opérationnelle, de diagnostic et de qualité de service.

Pour la CNAF, je peux apporter une capacité à tenir un rôle de \textbf{SPOC technique} avec méthode et calme : récupérer les contraintes du demandeur, faire émerger les dépendances, consolider les informations auprès des équipes techniques, identifier les points de vigilance, suivre les actions et produire un reporting clair. Je sais travailler avec des interlocuteurs variés, poser les bonnes questions, vulgariser les sujets complexes et garder une communication constructive lorsque les situations deviennent critiques.

Je dispose également d'un socle technique utile pour dialoguer avec les équipes d'architecture, build, run, production et support : Python, TypeScript, Angular, Node.js, Linux, Git, Docker, CI/CD, monitoring/logs, OpenSearch, Sentry, Confluence, Jira et Zenhub. Je ne me présente pas comme expert infrastructure DSI pur, mais comme un profil capable de comprendre les enjeux techniques, de les qualifier correctement et de faciliter leur traitement dans un cadre projet structuré.

La mission de la CNAF donne aussi du sens à cette candidature. Contribuer à une plateforme technique au service des équipes de développement, dans une DSI qui porte des services publics essentiels, correspond à ce que je recherche aujourd'hui : un poste utile, transverse, exigeant, où la rigueur technique et la qualité de coordination ont un impact concret.

Je serais heureux d'échanger avec vous afin de préciser comment mon expérience en coordination technique, validation, support run, documentation et analyse de risques peut contribuer aux missions de la Direction Technique.

Je vous prie d'agréer, Madame, Monsieur, l'expression de mes salutations distinguées.

\textbf{Guillaume Boileau}

}

\end{document}
